\honest{Outside the Laboratory}{Eliezer Yudkowsky}{2007}

``Outside the laboratory, scientists are no wiser than anyone else.'' Scientists say it humbly, and others say it to dismiss unwanted advice. Is the proverb true? ``Probably not in an absolute sense,'' but it ``does appear true to some degree,'' and that should alarm us. I give no evidence for the second part. \nb{In 2009 a Pew survey found American scientists about half as likely as the general public to believe in God or a higher power: 51\% against 95\%. That fits both the premise and its hedge.}

A shepherd trained to count sheep who stares blankly when asked to count apples has not understood counting. An economist who buys a lottery ticket every week may not understand expected utility. \nb{Whether the ticket is a mistake depends on how much the anticipation is worth, which is what the defenders in ``Lotteries: A Waste of Hope'' argue.} Feynman's students in Brazil had memorized optics: ask about Brewster's angle and they answered, but say ``Look at the water'' and nothing happened.

A competent scientist knows the routine: randomize treatment among N subjects, blind the judges, test for significance at 0.05. This is not etiquette like using the right fork; it is a way of testing hypotheses. Why do experiments? Not because journals demand them, or teachers taught it, or colleagues would look at you funny, but because to map a territory you must go and look; you cannot map a city from your living room with your eyes closed. Even checking your shoelaces needs photons from the Sun to bounce off them, strike your retina and be reconstructed by your visual cortex. Some physical process must correlate your brain with the world; reasoning is not magic.

Now take a scientist who believes in a spirit world and says, ``no one really knows, and I admit that I don't have any evidence''; it is a religious belief that observation cannot settle. ``I cannot but conclude that this person literally doesn't know why you have to look at things.'' If the scientist says the spirits spoke in their heart, that is a causal interaction and counts as evidence, to be weighed by its likelihood ratio against other causes of ``spirit voices'' and against the low prior of a complex belief with many parts; missing this is like a student not seeing that ``a medium with an index'' means water. People in lab coats say ``causal interaction'' and people in gaudy jewelry say ``spirits speaking,'' as if different clothing marked different realms; Gould's ``separate magisteria'' is an ``immortal blunder of a phrase.'' ``Causal interaction'' just means something that makes something else happen, and probability theory does not care what you wear.

To me, a scientist who accepts that spiritual matters cannot be settled by observation understands experiment only as a social convention: they know when experiments are expected, and where it is customary to make up beliefs without looking, they happily do that. A shepherd taught that seven sheep and eight make fifteen, or no dinner, might think seven apples and eight make three. If you know why the rules work, addition is the same for sheep and apples, and Newton is revered for showing that the planets obey the same laws as falling apples. In the everyday world different trees bear different fruit and different customs hold for different people, so a unified universe under fixed laws is deeply counterintuitive; only scientists really believe it. Feynman put it in a glass of wine: look closely and you see physics, the Earth's rocks, the age of the universe, the chemistry Pasteur traced to disease; our small minds divide it into sciences, but ``Nature does not know it!'' Some religions say everything is connected, which tells you no more than saying nothing is, and in practice make up one rule for girls under twelve, another for men over thirteen, one for the Sabbath, one for weekdays. Reality is one process under simple laws; different campus buildings are not different universes, and mind and matter, life and nonlife, are not divided; ``Nor is Bayes's Theorem different from one place to another.''

So a scientist who is as open to ``wacky ideas'' as anyone else outside their field ``probably never did understand why the scientific rules work,'' however well they parrot Popper. They ``don't like to be constrained by evidence''; they take off the lab coat at home and relax with comfortable nonsense, and it makes me wonder whether to trust them even in their own field. \nb{The scientist behind this diagnosis is the imagined one of the earlier paragraphs. The post gives no evidence about how real scientists who hold religious beliefs reason, inside or outside their fields.}

Maybe we can do better, if we learn enough probability theory to know why the rules work and enough psychology to see how they apply, if we learn to look at the water. That ambition gives up the comfortable modesty of being no better than anyone outside your specialty, but if theories of rationality do not apply to everyday life, something is wrong: it is not a different universe inside and outside the laboratory. ``A witty saying proves nothing,'' as Voltaire said.
